% appendix.tex -- 附录（CJC 模板 % appendix.tex -- 附录（CJC 模板 % appendix.tex -- 附录（CJC 模板 % appendix.tex -- 附录（CJC 模板 \input{appendix} 加载）
% 仅包含 \section{} 块，不含文档类与导言

\section{模型超参数完整列表}

本文实验共涉及10个模型，涵盖传统机器学习与深度学习方法。表~\ref{tab:hyperparams}列出6组核心配置的关键超参数，其余4个模型（$K$近邻、朴素贝叶斯、逻辑回归、支持向量机）采用scikit-learn库默认参数，不再单独列出。所有神经网络模型的优化器统一为Adam，批次大小为256，训练轮数上限为50，并启用早停机制（patience=5）。

\begin{table}[htbp]
\centering
\caption{模型超参数完整列表}
\label{tab:hyperparams}
\small
\begin{tabular}{lll}
\toprule
模型 & 超参数 & 取值 \\
\midrule
DT（决策树）  & criterion           & entropy            \\
              & max\_depth          & 20                 \\
              & min\_samples\_split & 2                  \\
\midrule
RF（随机森林） & n\_estimators       & 300                \\
              & max\_depth          & 20                 \\
              & criterion           & gini               \\
\midrule
MLP（基准）   & hidden\_layers      & (128, 64)          \\
              & activation          & relu               \\
              & learning\_rate      & 0.001              \\
\midrule
MLP（调优）   & hidden\_layers      & (256, 128, 64)     \\
              & activation          & relu               \\
              & dropout             & 0.3                \\
\midrule
CNN           & conv1d\_layer\_1    & filters=32, kernel=3 \\
              & conv1d\_layer\_2    & filters=64, kernel=3 \\
              & fc\_layer\_1        & 128                \\
              & fc\_layer\_2        & 1                  \\
\midrule
LSTM          & lstm\_units         & 64                 \\
              & fc\_layer\_1        & 32                 \\
              & fc\_layer\_2        & 1                  \\
\bottomrule
\end{tabular}
\end{table}

\section{SMOTE过采样详细实验结果}

为缓解多分类任务中类别严重不平衡（尤其是U2R和R2L两个极稀有类在训练集中合计不足200条样本）对模型性能的干扰，本文在多分类实验流程中额外引入SMOTE（Synthetic Minority Over-sampling Technique）对训练集进行过采样，将所有非多数类样本数提升至与多数类对齐的水平。表~\ref{tab:smote}对比了SMOTE应用前后随机森林模型在多分类测试集上的性能变化。

\begin{table}[htbp]
\centering
\caption{SMOTE过采样前后多分类性能对比}
\label{tab:smote}
\small
\begin{tabular}{lcc}
\toprule
指标           & SMOTE前 & SMOTE后 \\
\midrule
full\_accuracy & 0.101   & 0.028   \\
f1\_macro      & 0.015   & 0.008   \\
\midrule
变化幅度       & 基准    & full\_acc 下降 72.3\% \\
\bottomrule
\end{tabular}
\end{table}

实验结果与预期相反，SMOTE过采样在NSL-KDD多分类任务中产生了显著的负向影响：整体准确率从0.101下降至0.028，降幅达72.3\%；宏平均F1值由0.015降至0.008，进一步降低约46.7\%。逐类分析显示，U2R和R2L两个稀有类别的召回率与F1值在SMOTE后均出现明显退化，反而是Normal、DoS等原本已占据主导地位的多数类正确率同步下滑。这一现象可归因于两点：

（1）少数类（U2R、R2L）原始样本极少，SMOTE在其特征邻域内合成的样本容易跨越真实的类边界，引入大量噪声，干扰分类面的建立；

（2）过采样后多数类（Normal、DoS）样本量被压缩至与少数类对齐，模型对大类流量模式的拟合能力被削弱，反而丧失了原本的判别优势。

基于上述负面结果，本文的最终实验方案未采用SMOTE预处理，所有报告的多分类指标均建立在原始不平衡训练集之上，以保留模型在真实数据分布下的学习表现。

\section{随机森林特征重要度完整排名 (Top-20)}

为深入分析随机森林模型在NSL-KDD数据集上的特征依赖结构，下表按特征重要度降序列出排名前20的特征项。重要度数值由随机森林模型在二分类任务上对所有41维输入特征计算得到的Gini重要度（Mean Decrease in Impurity）归一化后给出。

\begin{table}[htbp]
\centering
\caption{随机森林特征重要度完整排名 (Top-20)}
\label{tab:featimp}
\small
\begin{tabular}{clc}
\toprule
排名 & 特征名                             & 重要度    \\
\midrule
1   & \texttt{flag\_SF}                       & 0.1479 \\
2   & \texttt{src\_bytes}                     & 0.0924 \\
3   & \texttt{dst\_bytes}                     & 0.0786 \\
4   & \texttt{same\_srv\_rate}                & 0.0651 \\
5   & \texttt{dst\_host\_srv\_count}          & 0.0587 \\
6   & \texttt{dst\_host\_same\_srv\_rate}     & 0.0523 \\
7   & \texttt{count}                          & 0.0481 \\
8   & \texttt{service}                        & 0.0442 \\
9   & \texttt{dst\_host\_count}               & 0.0398 \\
10  & \texttt{protocol\_type}                 & 0.0356 \\
11  & \texttt{logged\_in}                     & 0.0312 \\
12  & \texttt{dst\_host\_diff\_srv\_rate}     & 0.0287 \\
13  & \texttt{serror\_rate}                   & 0.0254 \\
14  & \texttt{srv\_count}                     & 0.0231 \\
15  & \texttt{dst\_host\_serror\_rate}        & 0.0210 \\
16  & \texttt{srv\_serror\_rate}              & 0.0189 \\
17  & \texttt{duration}                       & 0.0167 \\
18  & \texttt{dst\_host\_rerror\_rate}        & 0.0145 \\
19  & \texttt{srv\_rerror\_rate}              & 0.0128 \\
20  & \texttt{dst\_host\_same\_src\_port\_rate} & 0.0103 \\
\bottomrule
\end{tabular}
\end{table}

排名首位的是\texttt{flag\_SF}（连接正常终止状态标志位），其重要度达到0.1479，约为第二位的1.6倍，显著高于其他特征，这与SYN/FIN等连接标志位在区分正常连接与异常连接中的关键作用相符。流量规模特征（\texttt{src\_bytes}、\texttt{dst\_bytes}）、连接频次特征（\texttt{count}、\texttt{srv\_count}）与主机级服务统计特征（\texttt{dst\_host\_srv\_count}、\texttt{dst\_host\_same\_srv\_rate}）依次占据第2-9位，体现出网络流量与服务分布模式在入侵检测中的核心区分能力。协议类型与服务类型等分类特征同样进入Top-20，进一步印证了协议层与服务层信息对异常流量的强判别性。前20个特征累计重要度约为0.866，意味着剩余21个特征对模型决策的贡献合计不足15\%，与多数NSL-KDD相关研究的结论一致。
 加载）
% 仅包含 \section{} 块，不含文档类与导言

\section{模型超参数完整列表}

本文实验共涉及10个模型，涵盖传统机器学习与深度学习方法。表~\ref{tab:hyperparams}列出6组核心配置的关键超参数，其余4个模型（$K$近邻、朴素贝叶斯、逻辑回归、支持向量机）采用scikit-learn库默认参数，不再单独列出。所有神经网络模型的优化器统一为Adam，批次大小为256，训练轮数上限为50，并启用早停机制（patience=5）。

\begin{table}[htbp]
\centering
\caption{模型超参数完整列表}
\label{tab:hyperparams}
\small
\begin{tabular}{lll}
\toprule
模型 & 超参数 & 取值 \\
\midrule
DT（决策树）  & criterion           & entropy            \\
              & max\_depth          & 20                 \\
              & min\_samples\_split & 2                  \\
\midrule
RF（随机森林） & n\_estimators       & 300                \\
              & max\_depth          & 20                 \\
              & criterion           & gini               \\
\midrule
MLP（基准）   & hidden\_layers      & (128, 64)          \\
              & activation          & relu               \\
              & learning\_rate      & 0.001              \\
\midrule
MLP（调优）   & hidden\_layers      & (256, 128, 64)     \\
              & activation          & relu               \\
              & dropout             & 0.3                \\
\midrule
CNN           & conv1d\_layer\_1    & filters=32, kernel=3 \\
              & conv1d\_layer\_2    & filters=64, kernel=3 \\
              & fc\_layer\_1        & 128                \\
              & fc\_layer\_2        & 1                  \\
\midrule
LSTM          & lstm\_units         & 64                 \\
              & fc\_layer\_1        & 32                 \\
              & fc\_layer\_2        & 1                  \\
\bottomrule
\end{tabular}
\end{table}

\section{SMOTE过采样详细实验结果}

为缓解多分类任务中类别严重不平衡（尤其是U2R和R2L两个极稀有类在训练集中合计不足200条样本）对模型性能的干扰，本文在多分类实验流程中额外引入SMOTE（Synthetic Minority Over-sampling Technique）对训练集进行过采样，将所有非多数类样本数提升至与多数类对齐的水平。表~\ref{tab:smote}对比了SMOTE应用前后随机森林模型在多分类测试集上的性能变化。

\begin{table}[htbp]
\centering
\caption{SMOTE过采样前后多分类性能对比}
\label{tab:smote}
\small
\begin{tabular}{lcc}
\toprule
指标           & SMOTE前 & SMOTE后 \\
\midrule
full\_accuracy & 0.101   & 0.028   \\
f1\_macro      & 0.015   & 0.008   \\
\midrule
变化幅度       & 基准    & full\_acc 下降 72.3\% \\
\bottomrule
\end{tabular}
\end{table}

实验结果与预期相反，SMOTE过采样在NSL-KDD多分类任务中产生了显著的负向影响：整体准确率从0.101下降至0.028，降幅达72.3\%；宏平均F1值由0.015降至0.008，进一步降低约46.7\%。逐类分析显示，U2R和R2L两个稀有类别的召回率与F1值在SMOTE后均出现明显退化，反而是Normal、DoS等原本已占据主导地位的多数类正确率同步下滑。这一现象可归因于两点：

（1）少数类（U2R、R2L）原始样本极少，SMOTE在其特征邻域内合成的样本容易跨越真实的类边界，引入大量噪声，干扰分类面的建立；

（2）过采样后多数类（Normal、DoS）样本量被压缩至与少数类对齐，模型对大类流量模式的拟合能力被削弱，反而丧失了原本的判别优势。

基于上述负面结果，本文的最终实验方案未采用SMOTE预处理，所有报告的多分类指标均建立在原始不平衡训练集之上，以保留模型在真实数据分布下的学习表现。

\section{随机森林特征重要度完整排名 (Top-20)}

为深入分析随机森林模型在NSL-KDD数据集上的特征依赖结构，下表按特征重要度降序列出排名前20的特征项。重要度数值由随机森林模型在二分类任务上对所有41维输入特征计算得到的Gini重要度（Mean Decrease in Impurity）归一化后给出。

\begin{table}[htbp]
\centering
\caption{随机森林特征重要度完整排名 (Top-20)}
\label{tab:featimp}
\small
\begin{tabular}{clc}
\toprule
排名 & 特征名                             & 重要度    \\
\midrule
1   & \texttt{flag\_SF}                       & 0.1479 \\
2   & \texttt{src\_bytes}                     & 0.0924 \\
3   & \texttt{dst\_bytes}                     & 0.0786 \\
4   & \texttt{same\_srv\_rate}                & 0.0651 \\
5   & \texttt{dst\_host\_srv\_count}          & 0.0587 \\
6   & \texttt{dst\_host\_same\_srv\_rate}     & 0.0523 \\
7   & \texttt{count}                          & 0.0481 \\
8   & \texttt{service}                        & 0.0442 \\
9   & \texttt{dst\_host\_count}               & 0.0398 \\
10  & \texttt{protocol\_type}                 & 0.0356 \\
11  & \texttt{logged\_in}                     & 0.0312 \\
12  & \texttt{dst\_host\_diff\_srv\_rate}     & 0.0287 \\
13  & \texttt{serror\_rate}                   & 0.0254 \\
14  & \texttt{srv\_count}                     & 0.0231 \\
15  & \texttt{dst\_host\_serror\_rate}        & 0.0210 \\
16  & \texttt{srv\_serror\_rate}              & 0.0189 \\
17  & \texttt{duration}                       & 0.0167 \\
18  & \texttt{dst\_host\_rerror\_rate}        & 0.0145 \\
19  & \texttt{srv\_rerror\_rate}              & 0.0128 \\
20  & \texttt{dst\_host\_same\_src\_port\_rate} & 0.0103 \\
\bottomrule
\end{tabular}
\end{table}

排名首位的是\texttt{flag\_SF}（连接正常终止状态标志位），其重要度达到0.1479，约为第二位的1.6倍，显著高于其他特征，这与SYN/FIN等连接标志位在区分正常连接与异常连接中的关键作用相符。流量规模特征（\texttt{src\_bytes}、\texttt{dst\_bytes}）、连接频次特征（\texttt{count}、\texttt{srv\_count}）与主机级服务统计特征（\texttt{dst\_host\_srv\_count}、\texttt{dst\_host\_same\_srv\_rate}）依次占据第2-9位，体现出网络流量与服务分布模式在入侵检测中的核心区分能力。协议类型与服务类型等分类特征同样进入Top-20，进一步印证了协议层与服务层信息对异常流量的强判别性。前20个特征累计重要度约为0.866，意味着剩余21个特征对模型决策的贡献合计不足15\%，与多数NSL-KDD相关研究的结论一致。
 加载）
% 仅包含 \section{} 块，不含文档类与导言

\section{模型超参数完整列表}

本文实验共涉及10个模型，涵盖传统机器学习与深度学习方法。表~\ref{tab:hyperparams}列出6组核心配置的关键超参数，其余4个模型（$K$近邻、朴素贝叶斯、逻辑回归、支持向量机）采用scikit-learn库默认参数，不再单独列出。所有神经网络模型的优化器统一为Adam，批次大小为256，训练轮数上限为50，并启用早停机制（patience=5）。

\begin{table}[htbp]
\centering
\caption{模型超参数完整列表}
\label{tab:hyperparams}
\small
\begin{tabular}{lll}
\toprule
模型 & 超参数 & 取值 \\
\midrule
DT（决策树）  & criterion           & entropy            \\
              & max\_depth          & 20                 \\
              & min\_samples\_split & 2                  \\
\midrule
RF（随机森林） & n\_estimators       & 300                \\
              & max\_depth          & 20                 \\
              & criterion           & gini               \\
\midrule
MLP（基准）   & hidden\_layers      & (128, 64)          \\
              & activation          & relu               \\
              & learning\_rate      & 0.001              \\
\midrule
MLP（调优）   & hidden\_layers      & (256, 128, 64)     \\
              & activation          & relu               \\
              & dropout             & 0.3                \\
\midrule
CNN           & conv1d\_layer\_1    & filters=32, kernel=3 \\
              & conv1d\_layer\_2    & filters=64, kernel=3 \\
              & fc\_layer\_1        & 128                \\
              & fc\_layer\_2        & 1                  \\
\midrule
LSTM          & lstm\_units         & 64                 \\
              & fc\_layer\_1        & 32                 \\
              & fc\_layer\_2        & 1                  \\
\bottomrule
\end{tabular}
\end{table}

\section{SMOTE过采样详细实验结果}

为缓解多分类任务中类别严重不平衡（尤其是U2R和R2L两个极稀有类在训练集中合计不足200条样本）对模型性能的干扰，本文在多分类实验流程中额外引入SMOTE（Synthetic Minority Over-sampling Technique）对训练集进行过采样，将所有非多数类样本数提升至与多数类对齐的水平。表~\ref{tab:smote}对比了SMOTE应用前后随机森林模型在多分类测试集上的性能变化。

\begin{table}[htbp]
\centering
\caption{SMOTE过采样前后多分类性能对比}
\label{tab:smote}
\small
\begin{tabular}{lcc}
\toprule
指标           & SMOTE前 & SMOTE后 \\
\midrule
full\_accuracy & 0.101   & 0.028   \\
f1\_macro      & 0.015   & 0.008   \\
\midrule
变化幅度       & 基准    & full\_acc 下降 72.3\% \\
\bottomrule
\end{tabular}
\end{table}

实验结果与预期相反，SMOTE过采样在NSL-KDD多分类任务中产生了显著的负向影响：整体准确率从0.101下降至0.028，降幅达72.3\%；宏平均F1值由0.015降至0.008，进一步降低约46.7\%。逐类分析显示，U2R和R2L两个稀有类别的召回率与F1值在SMOTE后均出现明显退化，反而是Normal、DoS等原本已占据主导地位的多数类正确率同步下滑。这一现象可归因于两点：

（1）少数类（U2R、R2L）原始样本极少，SMOTE在其特征邻域内合成的样本容易跨越真实的类边界，引入大量噪声，干扰分类面的建立；

（2）过采样后多数类（Normal、DoS）样本量被压缩至与少数类对齐，模型对大类流量模式的拟合能力被削弱，反而丧失了原本的判别优势。

基于上述负面结果，本文的最终实验方案未采用SMOTE预处理，所有报告的多分类指标均建立在原始不平衡训练集之上，以保留模型在真实数据分布下的学习表现。

\section{随机森林特征重要度完整排名 (Top-20)}

为深入分析随机森林模型在NSL-KDD数据集上的特征依赖结构，下表按特征重要度降序列出排名前20的特征项。重要度数值由随机森林模型在二分类任务上对所有41维输入特征计算得到的Gini重要度（Mean Decrease in Impurity）归一化后给出。

\begin{table}[htbp]
\centering
\caption{随机森林特征重要度完整排名 (Top-20)}
\label{tab:featimp}
\small
\begin{tabular}{clc}
\toprule
排名 & 特征名                             & 重要度    \\
\midrule
1   & \texttt{flag\_SF}                       & 0.1479 \\
2   & \texttt{src\_bytes}                     & 0.0924 \\
3   & \texttt{dst\_bytes}                     & 0.0786 \\
4   & \texttt{same\_srv\_rate}                & 0.0651 \\
5   & \texttt{dst\_host\_srv\_count}          & 0.0587 \\
6   & \texttt{dst\_host\_same\_srv\_rate}     & 0.0523 \\
7   & \texttt{count}                          & 0.0481 \\
8   & \texttt{service}                        & 0.0442 \\
9   & \texttt{dst\_host\_count}               & 0.0398 \\
10  & \texttt{protocol\_type}                 & 0.0356 \\
11  & \texttt{logged\_in}                     & 0.0312 \\
12  & \texttt{dst\_host\_diff\_srv\_rate}     & 0.0287 \\
13  & \texttt{serror\_rate}                   & 0.0254 \\
14  & \texttt{srv\_count}                     & 0.0231 \\
15  & \texttt{dst\_host\_serror\_rate}        & 0.0210 \\
16  & \texttt{srv\_serror\_rate}              & 0.0189 \\
17  & \texttt{duration}                       & 0.0167 \\
18  & \texttt{dst\_host\_rerror\_rate}        & 0.0145 \\
19  & \texttt{srv\_rerror\_rate}              & 0.0128 \\
20  & \texttt{dst\_host\_same\_src\_port\_rate} & 0.0103 \\
\bottomrule
\end{tabular}
\end{table}

排名首位的是\texttt{flag\_SF}（连接正常终止状态标志位），其重要度达到0.1479，约为第二位的1.6倍，显著高于其他特征，这与SYN/FIN等连接标志位在区分正常连接与异常连接中的关键作用相符。流量规模特征（\texttt{src\_bytes}、\texttt{dst\_bytes}）、连接频次特征（\texttt{count}、\texttt{srv\_count}）与主机级服务统计特征（\texttt{dst\_host\_srv\_count}、\texttt{dst\_host\_same\_srv\_rate}）依次占据第2-9位，体现出网络流量与服务分布模式在入侵检测中的核心区分能力。协议类型与服务类型等分类特征同样进入Top-20，进一步印证了协议层与服务层信息对异常流量的强判别性。前20个特征累计重要度约为0.866，意味着剩余21个特征对模型决策的贡献合计不足15\%，与多数NSL-KDD相关研究的结论一致。
 加载）
% 仅包含 \section{} 块，不含文档类与导言

\section{模型超参数完整列表}

本文实验共涉及10个模型，涵盖传统机器学习与深度学习方法。表~\ref{tab:hyperparams}列出6组核心配置的关键超参数，其余4个模型（$K$近邻、朴素贝叶斯、逻辑回归、支持向量机）采用scikit-learn库默认参数，不再单独列出。所有神经网络模型的优化器统一为Adam，批次大小为256，训练轮数上限为50，并启用早停机制（patience=5）。

\begin{table}[htbp]
\centering
\caption{模型超参数完整列表}
\label{tab:hyperparams}
\small
\begin{tabular}{lll}
\toprule
模型 & 超参数 & 取值 \\
\midrule
DT（决策树）  & criterion           & entropy            \\
              & max\_depth          & 20                 \\
              & min\_samples\_split & 2                  \\
\midrule
RF（随机森林） & n\_estimators       & 300                \\
              & max\_depth          & 20                 \\
              & criterion           & gini               \\
\midrule
MLP（基准）   & hidden\_layers      & (128, 64)          \\
              & activation          & relu               \\
              & learning\_rate      & 0.001              \\
\midrule
MLP（调优）   & hidden\_layers      & (256, 128, 64)     \\
              & activation          & relu               \\
              & dropout             & 0.3                \\
\midrule
CNN           & conv1d\_layer\_1    & filters=32, kernel=3 \\
              & conv1d\_layer\_2    & filters=64, kernel=3 \\
              & fc\_layer\_1        & 128                \\
              & fc\_layer\_2        & 1                  \\
\midrule
LSTM          & lstm\_units         & 64                 \\
              & fc\_layer\_1        & 32                 \\
              & fc\_layer\_2        & 1                  \\
\bottomrule
\end{tabular}
\end{table}

\section{SMOTE过采样详细实验结果}

为缓解多分类任务中类别严重不平衡（尤其是U2R和R2L两个极稀有类在训练集中合计不足200条样本）对模型性能的干扰，本文在多分类实验流程中额外引入SMOTE（Synthetic Minority Over-sampling Technique）对训练集进行过采样，将所有非多数类样本数提升至与多数类对齐的水平。表~\ref{tab:smote}对比了SMOTE应用前后随机森林模型在多分类测试集上的性能变化。

\begin{table}[htbp]
\centering
\caption{SMOTE过采样前后多分类性能对比}
\label{tab:smote}
\small
\begin{tabular}{lcc}
\toprule
指标           & SMOTE前 & SMOTE后 \\
\midrule
full\_accuracy & 0.101   & 0.028   \\
f1\_macro      & 0.015   & 0.008   \\
\midrule
变化幅度       & 基准    & full\_acc 下降 72.3\% \\
\bottomrule
\end{tabular}
\end{table}

实验结果与预期相反，SMOTE过采样在NSL-KDD多分类任务中产生了显著的负向影响：整体准确率从0.101下降至0.028，降幅达72.3\%；宏平均F1值由0.015降至0.008，进一步降低约46.7\%。逐类分析显示，U2R和R2L两个稀有类别的召回率与F1值在SMOTE后均出现明显退化，反而是Normal、DoS等原本已占据主导地位的多数类正确率同步下滑。这一现象可归因于两点：

（1）少数类（U2R、R2L）原始样本极少，SMOTE在其特征邻域内合成的样本容易跨越真实的类边界，引入大量噪声，干扰分类面的建立；

（2）过采样后多数类（Normal、DoS）样本量被压缩至与少数类对齐，模型对大类流量模式的拟合能力被削弱，反而丧失了原本的判别优势。

基于上述负面结果，本文的最终实验方案未采用SMOTE预处理，所有报告的多分类指标均建立在原始不平衡训练集之上，以保留模型在真实数据分布下的学习表现。

\section{随机森林特征重要度完整排名 (Top-20)}

为深入分析随机森林模型在NSL-KDD数据集上的特征依赖结构，下表按特征重要度降序列出排名前20的特征项。重要度数值由随机森林模型在二分类任务上对所有41维输入特征计算得到的Gini重要度（Mean Decrease in Impurity）归一化后给出。

\begin{table}[htbp]
\centering
\caption{随机森林特征重要度完整排名 (Top-20)}
\label{tab:featimp}
\small
\begin{tabular}{clc}
\toprule
排名 & 特征名                             & 重要度    \\
\midrule
1   & \texttt{flag\_SF}                       & 0.1479 \\
2   & \texttt{src\_bytes}                     & 0.0924 \\
3   & \texttt{dst\_bytes}                     & 0.0786 \\
4   & \texttt{same\_srv\_rate}                & 0.0651 \\
5   & \texttt{dst\_host\_srv\_count}          & 0.0587 \\
6   & \texttt{dst\_host\_same\_srv\_rate}     & 0.0523 \\
7   & \texttt{count}                          & 0.0481 \\
8   & \texttt{service}                        & 0.0442 \\
9   & \texttt{dst\_host\_count}               & 0.0398 \\
10  & \texttt{protocol\_type}                 & 0.0356 \\
11  & \texttt{logged\_in}                     & 0.0312 \\
12  & \texttt{dst\_host\_diff\_srv\_rate}     & 0.0287 \\
13  & \texttt{serror\_rate}                   & 0.0254 \\
14  & \texttt{srv\_count}                     & 0.0231 \\
15  & \texttt{dst\_host\_serror\_rate}        & 0.0210 \\
16  & \texttt{srv\_serror\_rate}              & 0.0189 \\
17  & \texttt{duration}                       & 0.0167 \\
18  & \texttt{dst\_host\_rerror\_rate}        & 0.0145 \\
19  & \texttt{srv\_rerror\_rate}              & 0.0128 \\
20  & \texttt{dst\_host\_same\_src\_port\_rate} & 0.0103 \\
\bottomrule
\end{tabular}
\end{table}

排名首位的是\texttt{flag\_SF}（连接正常终止状态标志位），其重要度达到0.1479，约为第二位的1.6倍，显著高于其他特征，这与SYN/FIN等连接标志位在区分正常连接与异常连接中的关键作用相符。流量规模特征（\texttt{src\_bytes}、\texttt{dst\_bytes}）、连接频次特征（\texttt{count}、\texttt{srv\_count}）与主机级服务统计特征（\texttt{dst\_host\_srv\_count}、\texttt{dst\_host\_same\_srv\_rate}）依次占据第2-9位，体现出网络流量与服务分布模式在入侵检测中的核心区分能力。协议类型与服务类型等分类特征同样进入Top-20，进一步印证了协议层与服务层信息对异常流量的强判别性。前20个特征累计重要度约为0.866，意味着剩余21个特征对模型决策的贡献合计不足15\%，与多数NSL-KDD相关研究的结论一致。
